% !TEX root = report.tex

\section{Related Work}
\label{sec:related_work}

\textbf{Wi-Fi sensing and the SHARP system.}
Device-free HAR from commodity Wi-Fi has matured rapidly thanks to the
availability of CSI extraction tools and public datasets. The most relevant work
for us is SHARP~\cite{sharp}, which recognizes human activities from the Doppler
information contained in the CSI of IEEE~802.11ac links and, crucially, targets
\emph{environment and person independence}. SHARP achieves this through two
ingredients: (i) a robust phase-sanitization step that isolates the multipath
components and adopts the \emph{strongest path} as the phase reference, and (ii)
a \emph{decision-fusion} scheme that runs a compact Inception-style CNN on each
of the four receiver antennas and combines the per-antenna decisions. On the
seven measurement sets of the accompanying CSI dataset~\cite{csidataset}, the
full system exceeds $95\%$ accuracy even on rooms, days and persons never seen at
training time. These results are excellent, but they rest on a heavy,
hand-engineered signal-processing front-end; the paper's own ablation shows that
replacing the strongest-path sanitization with the standard reference-antenna
method makes the \emph{same} CNN collapse to $39$--$60\%$ on the unseen sets,
revealing how much of the robustness is carried by the pre-processing rather than
by the learned representation. Our work keeps the minimal front-end (single
antenna, no strongest-path sanitization) and asks how much cross-domain
robustness can instead be recovered on the \emph{representation} side.

\textbf{Contrastive representation learning.}
Contrastive methods learn representations by pulling together ``views'' that
should be similar and pushing apart those that should differ. SimCLR~\cite{simclr}
established a simple self-supervised recipe: two random augmentations of the same
image form a positive pair, all other images in the batch are negatives, and an
encoder is trained with a temperature-scaled InfoNCE loss on $L_2$-normalized
embeddings. Supervised Contrastive Learning (SupCon)~\cite{supcon} extends this
idea to the labelled setting: \emph{all} samples sharing the anchor's label
become positives, which produces tighter, more class-consistent clusters than
cross-entropy and improves robustness to natural corruptions. We adopt SupCon
because its notion of ``positive'' maps naturally onto our goal: two windows of
the same activity should be neighbours \emph{regardless of the environment,
person or antenna they come from}, so that the encoder is driven to discard
domain-specific nuisances. To the best of our knowledge, combining SupCon with
explicit domain-randomization augmentations for cross-domain Wi-Fi HAR on the
SHARP data is the specific extension we contribute here: it complements the
signal-processing route of~\cite{sharp} by attacking domain shift from the
feature-learning side, and we quantify both its gains and its residual limits in
\secref{sec:results}.
